\title{Direct Preference Optimization: Your Language Model is Secretly a Reward Model}

\begin{abstract}
Large-scale unsupervised language models (LMs) acquire extensive world knowledge and develop some reasoning abilities through their training; however, their training paradigm, being entirely unsupervised, makes fine-grained regulation of their outputs challenging. Common approaches to address this limitation rely on gathering human preference data—where annotators compare the quality of different model outputs—and subsequently fine-tuning the original LM to better adhere to these labeled preferences, most notably using reinforcement learning from human feedback (RLHF). Nevertheless, the RLHF framework typically involves a multi-step and potentially unstable workflow: it first trains a reward model to encapsulate the recorded human preferences, and subsequently applies reinforcement learning to adapt the LM so as to maximize this modeled reward, all the while maintaining proximity to the original distribution. In this work, we propose a novel reward model parameterization for RLHF that yields the optimal policy in closed form, thereby reducing the standard RLHF objective to a straightforward classification loss. We term our approach \textit{Direct Preference Optimization} (DPO), which demonstrates stability, efficiency, and modest computational demands—removing the necessity for LM sampling during the fine-tuning phase as well as reducing the extent of required hyperparameter tuning. Empirically, we find that DPO fine-tunes LMs to follow human preferences at least as well as, and in certain aspects better than, prevalent alternatives. In particular, DPO outperforms RLHF methods based on PPO for sentiment control, and achieves comparable or superior results for summarization and single-turn dialogue quality, all while offering considerable simplicity in both implementation and training.
\end{abstract}

\section{Introduction}
Large-scale unsupervised language models (LMs) trained on immense corpora have demonstrated remarkable emergent abilities~\citep{chowdhery2022palm, brown2020language, touvron2023llama,bubeck2023sparks}. Despite these achievements, such models are trained on data generated by humans whose objectives, values, and expertise can be highly diverse. Not all these human-derived tendencies are desirable; for instance, while an AI programming assistant should be able to \textit{recognize} common coding errors to help remediate them, we would generally prefer the model to \emph{generate} solutions that reflect the rare instances of expert coding proficiency found within its training set. Likewise, although awareness of widespread misconceptions—those held by a significant portion of people—might be useful, we do not want our model to repeat these falsehoods in half its answers merely because they are commonly believed. Thus, it is essential to \emph{select and enforce preferred responses and behaviors} from the model’s broad \textit{reservoir of knowledge and skills}, thereby ensuring that AI systems are both safe and responsive to user needs, as well as reliably controllable \citep{ouyang2022training}. Conventionally, methods for this behavioral alignment employ reinforcement learning (RL) to align LMs with human values, but as we will demonstrate, the RL-driven objective behind current approaches can be exactly solved by leveraging a straightforward binary cross-entropy loss—substantially streamlining the preference learning workflow.

In general, current alignment strategies aim to imbue language models with intended behaviors through carefully constructed datasets reflecting human preferences about desirable and helpful conduct. This alignment process, often referred to as preference learning, follows an initial period of extensive unsupervised pre-training on text corpora. The simplest implementation is supervised fine-tuning on manually created examples of high-quality output, but the prevailing state-of-the-art involves reinforcement learning from human (or AI) feedback—RLHF or RLAIF \citep{christiano2017deep,bai2022constitutional}. Here, a reward model is trained on annotated preference data, and then RL is used to steer the LM’s policy towards responses with higher reward, while regularizing against substantial divergence from its pre-trained version. Although RLHF has led to language models with strong capabilities in conversation and code generation, the underlying process is considerably more involved than standard supervised fine-tuning. It necessitates training multiple models and requires policy sampling throughout optimization, contributing to a significant computational overhead.

In this work, we demonstrate a method for optimizing language models to align with human preferences without relying on explicit reward modeling or reinforcement learning techniques. We introduce \textit{Direct Preference Optimization} (DPO), an approach that, while implicitly optimizing the same objective as established RLHF algorithms (maximizing reward under a KL-divergence regularization), offers a straightforward implementation and training procedure. Conceptually, DPO adjusts the relative log likelihood in favor of preferred over non-preferred outputs, yet uniquely applies a per-example, adaptive importance weighting mechanism to avoid the degeneration of the model that occurs under a naive application of the probability ratio objective. Similar to existing strategies, DPO employs a theoretical preference framework (for example, the Bradley-Terry model; \cite{bradley1952rankanalysis}) for evaluating the agreement between a reward function and actual human preference data. In contrast to conventional approaches, which utilize this preference model to construct a loss for training a reward function before fitting a policy to the learned reward, DPO leverages a change-of-variables technique to formulate the preference loss directly as a function of the policy. As a result, with access to a collection of human preference judgements over model outputs, DPO enables direct optimization of the policy using a binary cross entropy loss, yielding the optimal policy corresponding to an implicit reward model shaped by the preference data.

The primary contribution presented here is Direct Preference Optimization (DPO), a streamlined, reinforcement learning-free method for preference-based language model training. Empirical results indicate that DPO performs on par with, or better than, leading methods such as PPO-based RLHF in learning from human preferences across various tasks—including sentiment control, text summarization, and conversational response generation—applied to models containing up to 6 billion parameters.

\section{Related Work}

As self-supervised language models grow larger, they increasingly demonstrate the ability to solve some problems in a zero-shot manner \citep{radford2019language} or with only a handful of prompt examples \citep{gpt3,megatron,chowdhery2022palm}. Nonetheless, further improvements in performance on downstream applications and better alignment with user expectations are typically achieved by fine-tuning on corpora comprised of instructions paired with human-generated completions \citep{mishra-etal-2022-cross,sanh2022multitask,chung2022scaling,thoppilan2022lamda}. The process known as `instruction-tuning’ enhances the capacity of large language models (LLMs) to generalize to novel instructions not present in the fine-tuning data, boosting their overall effectiveness \citep{chung2022scaling}. Yet, although instruction tuning yields notable improvements, collecting \textit{relative} assessments of completion quality from humans is often less labor-intensive than gathering expert-authored exemplars. Accordingly, later research has prioritized fine-tuning LLMs using datasets that encode human preferences, leading to advancements in areas such as machine translation \citep{kreutzer-etal-2018-reliability}, text summarization \citep{stiennon2022learning,ziegler2020finetuning}, creative writing \citep{ziegler2020finetuning}, and adherence to user instructions \citep{ouyang2022training,ramamurthy2023is}. These techniques initially train a neural reward function to model the preference data—typically utilizing a preference-learning model like the Bradley-Terry model \citep{bradley1952rankanalysis}—and subsequently employ reinforcement learning approaches, such as REINFORCE \citep{williams1992reinforce}, proximal policy optimization (PPO; \cite{schulman2017proximal}), or related strategies \citep{ramamurthy2023is}, to fine-tune the language model for reward maximization. Related efforts use LLMs that are already instruction-following and preference-aligned to synthesize additional preference-annotated data, targeting characteristics like safety or non-harmfulness \citep{bai2022constitutional}, while relying on lightweight supervision such as text-based rubrics provided by humans to guide the annotations. Collectively, these advances integrate research from two central domains: reinforcement learning techniques applied to language models for various training objectives \citep{Ranzato2015SequenceLT,paulus2018a,wu2018learning} and broad strategies for learning from human preference feedback \citep{christiano2017deep,kupcsik2018learning}. Even with the promise shown by leveraging relative preference judgments, the reinforcement learning-based fine-tuning of large language models continues to present formidable practical hurdles; this study introduces a theoretically grounded methodology for optimizing over relative preferences that does not depend on reinforcement learning.

Independent of linguistic applications, the study of learning policies based on preferences has a history within both bandit and reinforcement learning paradigms, leading to the development of diverse strategies. The contextual dueling bandit (CDB; \cite{yue2012karmed,dudik2015contextual}) formalism, for instance, investigates learning within contextual bandits where feedback is delivered through preferences or rankings over actions, rather than explicit reward signals. When absolute rewards are unavailable, CDBs replace the standard concept of an optimal policy with that of a \textit{von Neumann winner}—a policy expected to defeat any alternative policy with probability at least 50\% \citep{dudik2015contextual}. Notably, CDBs assume access to preferences obtained sequentially (online); in contrast, preference learning from human data often involves training on a fixed dataset comprising pre-annotated action pairs with preference labels \citep{yan2022human}. Similarly, \textit{preference-based RL} (PbRL) operates by interpreting binary preference information supplied by an opaque 'scoring' function in lieu of reward feedback \citep{BusaFekete2014,ruiz2023dueling}. Several PbRL algorithms have been introduced, some capable of utilizing off-policy preference data. Nevertheless, these approaches typically require an initial phase to estimate the hidden scoring (reward) function, which is subsequently used for policy optimization \citep{jain2013learning,BusaFekete2014,christiano2017deep,sadigh2017active,kupcsik2018learning}. Our contribution departs from these methods by directly optimizing policies to align with observed preferences, thus consolidating the learning process into a single stage.

\section{Preliminaries}\label{section:prelims}

We summarize the RLHF pipeline described by \citeauthor{ziegler2020finetuning} and expanded in subsequent work (\citep{stiennon2022learning, bai2022training, ouyang2022training}). The pipeline is conventionally structured into three primary phases: (1) supervised fine-tuning (SFT); (2) preference data collection and reward function estimation; and (3) reinforcement learning for policy optimization.

\textbf{SFT}: The RLHF workflow commences by adapting a large pre-trained language model to downstream tasks of interest (such as dialogue or summarization) through supervised learning on curated datasets, yielding a policy denoted as $\pi^\text{SFT}$.

\textbf{Reward Modelling Phase}: Subsequently, the SFT model is conditioned on prompts $x$ to produce paired outputs $(y_1, y_2)\sim \pi^\text{SFT}(y \mid x)$. Human annotators are then asked to indicate their preference between the two responses, recorded as $y_w\succ y_l \mid x$, where $y_w$ is the favored and $y_l$ the less favored completion from the generated pair. It is assumed these preferences reflect evaluations according to an unobserved reward function $r^*(y, x)$. To model such preference data, frameworks like the Bradley-Terry (BT) model \cite{bradley1952rankanalysis} are commonly used, though the Plackett-Luce model \citep{plackett1975analysis, luce2012individual} can be employed for more complex ranked sets. According to the BT model, the likelihood of observed human preferences $p^*$ is given by:

Given a static dataset of comparisons $\mathcal{D}=\bigl\{x^{(i)}, y_w^{(i)}, y_l^{(i)}\bigr\}_{i=1}^N$ drawn from $p^*$, we can introduce a reward model $r_{\phi}(x, y)$ and estimate its parameters through maximum likelihood estimation. When the problem is posed as a binary classification task, we utilize the negative log-likelihood loss:

\begin{equation}\label{eq:reward_model}
    \mathcal{L}_R(r_{\phi}, \mathcal{D}) = -\mathbb{E}_{(x, y_w, y_l)\sim \mathcal{D}}\bigl[\log \sigma(r_{\phi}(x, y_w)- r_{\phi}(x, y_l))\bigr]
\end{equation}

with $\sigma$ denoting the logistic sigmoid function. In the context of language models, the reward network $r_{\phi}(x, y)$ is frequently initialized from the SFT model $\pi^\text{SFT}(y \mid x)$ and augmented with an additional linear output layer placed atop the final transformer layer, generating a single scalar reward \cite{ziegler2020finetuning}. To mitigate variance in the learned reward function, previous work has enforced reward normalization such that $\mathbb{E}_{x,y\sim \mathcal{D}}\left[r_\phi(x, y)\right] = 0$ holds for all $x$.

\textbf{RL Fine-Tuning Phase}: During this phase, the reward function that has been trained serves as a feedback mechanism for improving the language model. As outlined in previous studies~\citep{jaques2017sequence, jaques2020human}, the problem is formulated as:

\begin{equation}\label{eq:RL}
\max_{\pi_{\theta}}  \mathbb{E}_{x\sim \mathcal{D}, y\sim \pi_{\theta}(y \mid x)}\bigl[r_{\phi}(x, y)\bigr] - \beta\mathbb{D}_{\textrm{KL}}\bigl[\pi_{\theta}(y\mid x)\mid \mid \pi_\text{ref}(y\mid x)\bigr],
\end{equation}

where $\beta$ regulates the extent of divergence from the reference policy $\pi_\text{ref}$—typically the pretrained SFT model $\pi^\text{SFT}$. In actual implementations, the language model policy $\pi_\theta$ is also initialized with $\pi^\text{SFT}$. This KL-regularization is essential to ensure that the policy does not depart excessively from regions of the data distribution where the reward model is valid, maintaining output diversity and avoiding over-concentration on a small set of high-reward responses. Due to the non-differentiable characteristics of language outputs, this objective is generally addressed through reinforcement learning. The prevalent method \citep{ziegler2020finetuning, stiennon2022learning, bai2022training, ouyang2022training} constructs the reward as ${r(x, y) = r_{\phi}(x, y) -\beta (\log \pi_{\theta}(y\mid x) - \log \pi_\text{ref}(y\mid x))}$ and applies PPO \cite{schulman2017proximal} to maximize it.

\section{Direct Preference Optimization}\label{sec:DPO}

Seeking to bypass the intricacies of reinforcement learning at the scale necessary for large language model fine-tuning, we aim to develop a streamlined method for optimizing policies directly from preference data. Departing from traditional RLHF pipelines—which separate reward model training and RL-based policy optimization—our methodology takes advantage of a specific reward function parameterization that permits analytic recovery of the corresponding optimal policy, obviating the need for reinforcement learning iterations. We capitalize on an explicit, closed-form correspondence between the reward function and the optimal policy, allowing us to convert a loss on rewards into one on policies. This transformation eliminates the requirement for fitting a distinct reward model, while retaining the interpretive framework provided by established preference models such as the Bradley-Terry model. Accordingly, the policy network concurrently encapsulates the functionalities of both the language generator and the (implicit) reward estimator.

\textbf{Reformulating the DPO objective.} We begin with the RL objective from prior studies, denoted as Eq.~\ref{eq:RL}, under a generic reward function $r$. Prior works~\citep{peters2007reinforcement, peng2019advantage, korbak2022reinforcement, go2023aligning} have established that the solution to the reward maximization problem constrained by KL-divergence, as posed in Eq.~\ref{eq:RL}, yields the following optimal policy:

\begin{equation}\label{eq:op_policy}
    \pi_r(y\mid x) = \frac{1}{Z(x)}\pi_\text{ref}(y\mid x)\exp\left(\frac{1}{\beta}r(x, y)\right),
\end{equation}

where $Z(x) =\sum_{y}\pi_\text{ref}(y\mid x)\exp\left(\frac{1}{\beta}r(x, y)\right)$ serves as the normalization constant. A full derivation is provided in Appendix \ref{app:derivation1}. Although one may use the MLE estimate $r_{\phi}$ for the unknown ground-truth reward $r^*$, computing the partition function $Z(x)$ remains computationally demanding~\citep{korbak2022reinforcement, go2023aligning}, posing significant practical challenges. Nevertheless, Eq.~\ref{eq:op_policy} can be rearranged to express the reward function via its optimal policy $\pi_r$, the reference $\pi_\text{ref}$, and the partition function $Z(\cdot)$. Specifically, by applying a logarithm to both sides of Eq.~\ref{eq:op_policy} and simplifying, we arrive at:

\begin{equation}\label{eq:main_eq}
    r(x,y) =\beta \log \frac{\pi_r(y\mid x)}{\pi_\text{ref}(y\mid x)} + \beta \log Z(x).
\end{equation}

This form allows us to replace the true reward $r^*$ using its associated optimal policy $\pi^*$. Importantly, within the Bradley-Terry framework, outcomes depend solely on reward differences between two samples, namely, ${p^*(y_1 \succ y_2 \mid x) = \sigma(r^*(x, y_1) - r^*(x, y_2))}$. Upon substituting the parameterization from Eq.~\ref{eq:main_eq} into the preference equation (Eq.~\ref{eq:bradley-terry}), the terms involving the partition function eliminate, leading to a preference probability that depends exclusively on the optimal and reference policies. Thus, the optimal policy $\pi^*$ for RLHF under the Bradley-Terry model fulfills:

\begin{equation}\label{eq:objective}
    p^*(y_1\succ y_2 \mid x)=\frac{1}{1 + \exp\left(\beta \log \frac{\pi^*(y_2\mid x)}{\pi_\text{ref}(y_2\mid x)} - \beta \log \frac{\pi^*(y_1\mid x)}{\pi_\text{ref}(y_1\mid x)}\right)}
\end{equation}

Complete steps are outlined in Appendix~\ref{app:derivation2}. While Eq.~\ref{eq:objective} is constructed for the Bradley-Terry model, an analogous methodology extends to the broader Plackett-Luce models~\citep{plackett1975analysis, luce2012individual}, discussed further in Appendix~\ref{app:plackett_luce_models}.

By expressing the preference likelihood using the optimal policy (as opposed to a reward model), it is possible to formulate a maximum likelihood estimation objective for a parametrized policy $\pi_\theta$. Mirroring the reward-based approach (cf. Eq.~\ref{eq:reward_model}), our training objective for policies becomes:

\begin{equation}\label{eq:optimum_model}
    \mathcal{L}_\text{DPO}(\pi_{\theta}; \pi_\text{ref}) = -\mathbb{E}_{(x, y_w, y_l)\sim \mathcal{D}}\left[\log \sigma \left(\beta \log \frac{\pi_{\theta}(y_w\mid x)}{\pi_\text{ref}(y_w\mid x)} - \beta \log \frac{\pi_{\theta}(y_l\mid x)}{\pi_\text{ref}(y_l\mid x)}\right)\right].
\end{equation

In this manner, we model an implicit reward by adopting an alternative parameterization, ensuring that the resulting optimal policy coincides with $\pi_\theta$. Furthermore, because this approach is equivalent to learning a reparametrized Bradley-Terry model, it inherits several theoretical advantages, such as consistency under appropriate assumptions on the distribution of preference data \cite{bong2022generalized}. Additional theoretical aspects of DPO and its connections to related methodologies are elaborated in Section~\ref{sec:theory}.

\textbf{How does the DPO update function?} To gain mechanistic insight into DPO, it is instructive to examine the gradient of its loss function, $\mathcal{L}_\text{DPO}$. The parameter gradient can be expressed as:

\begin{multline*}\label{eq:gradient}
    \nabla_\theta \mathcal{L}_\text{DPO}(\pi_\theta;\pi_\text{ref}) = \\ -\beta\mathbb{E}_{(x, y_w, y_l) \sim \mathcal{D}} \bigg[\underbrace{\sigma(\hat{r}_\theta(x, y_l) - \hat{r}_\theta (x, y_w))}_\text{higher weight when reward estimate is wrong}\bigg[\underbrace{\nabla_\theta\log \pi(y_w \mid x)}_\text{increase likelihood of $y_w$} - \underbrace{\nabla_\theta\log\pi(y_l \mid x)}_\text{decrease likelihood of $y_l$}\bigg]\bigg],
\end{multline*}

Here, $\hat{r}_\theta(x, y) = \beta \log \frac{\pi_\theta(y \mid x)}{\pi_\text{ref}(y \mid x)}$ defines the implicit reward in terms of the language model $\pi_\theta$ and a reference model $\pi_\text{ref}$ (see Section~\ref{sec:theory} for further details). In essence, the loss gradient $\mathcal{L}_\text{DPO}$ operates by enhancing the likelihood of favorable completions $y_w$ and reducing the likelihood of less preferred completions $y_l$. Crucially, each example is assigned a weight according to how much the implicit reward model $\hat{r}_\theta$ erroneously favors $y_l$ over $y_w$, modulated by $\beta$. This reflects the extent to which the model incorrectly ranks the options, as well as the influence of the KL regularization. Empirical evidence from our studies underscores the necessity of this weighting; omitting it (as in a simplified variant of the algorithm) can result in catastrophic degradation of language model performance (see Appendix Table~\ref{tab:unlikelihood_generations}).

\textbf{DPO overview.} The standard DPO workflow consists of two main steps: (1) For each prompt $x$, generate completion samples $y_1, y_2 \sim \pi_\text{ref}(\cdot \mid x)$, then annotate these with human preference judgments to build an offline preference dataset $\mathcal{D} = \{x^{(i)}, y_w^{(i)}, y_l^{(i)}\}_{i=1}^N$; (2) update the language model $\pi_\theta$ by minimizing $\mathcal{L}_\text{DPO}$, using the chosen $\pi_\text{ref}$, dataset $\mathcal{D}$, and the designated $\beta$ parameter.

In actual implementation, it is preferable to make use of existing public preference datasets, rather than collecting new samples or soliciting additional human preferences. As these datasets are typically collected from completions produced by $\pi^\text{SFT}$, we initialize $\pi_\text{ref} = \pi^\text{SFT}$ wherever possible. If $\pi^\text{SFT}$ is unavailable, we instead fit $\pi_\text{ref}$ by maximizing the likelihood over the preferred completions ${(x, y_w)}$, specifically setting ${\pi_\text{ref} = \argmax_{\pi}\mathbb{E}_{x, y_w \sim \mathcal{D}}\left[\log \pi(y_w \mid x)\right]}$. This initialization strategy serves to lessen any distributional mismatch between the ideal—but inaccessibly defined—reference distribution and the $\pi_\text{ref}$ actually utilized during DPO training. Additional details on implementation and the choice of hyperparameters are provided in Appendix~\ref{app:implementation}.

\section{Theoretical Analysis of DPO} This section aims to offer a deeper conceptual understanding of DPO, present supporting theoretical results, and elucidate why DPO addresses certain challenges associated with actor-critic frameworks in RLHF, such as those encountered with PPO~\cite{schulman2017proximal}.

\label{sec:theory}

\subsection{Your Language Model Is Secretly a Reward Model} DPO enables simultaneous avoidance of explicit reward modeling and separate RL policy optimization by relying solely on a maximum likelihood framework. Importantly, the objective in Eq. \ref{eq:main_eq} is equivalent to a Bradley-Terry model with reward given by $r^*(x, y) = \beta \log\frac{\pi^*_\theta(y \mid x)}{\pi_\text{ref}(y \mid x)}$. Thus, optimizing the model $\pi_\theta$ amounts to optimizing the reward model as in Eq. \ref{eq:reward_model} under a reparameterization. The rest of this section builds theoretical foundations for this perspective, demonstrates that this reparameterization does not restrict the class of reward models that can be learned, and shows that it can recover the optimal policy exactly. We start by introducing an equivalence relation on reward functions.

\begin{definition}
We define two reward functions $r(x, y)$ and $r'(x, y)$ as equivalent if and only if ${r(x, y)-r'(x, y) = f(x)}$ for some function $f$.     
\end{definition}

It is straightforward to verify that this constitutes an equivalence relation, dividing the set of reward functions into distinct equivalence classes. This leads to two immediate lemmas:

\begin{lemma}\label{lemma:same_prefrence} In the context of the Plackett-Luce model—and, in particular, the Bradley-Terry preference framework—any pair of reward functions from the same equivalence class produces identical preference distributions.
\end{lemma}

\begin{lemma}\label{lemma:same_policy}
    For the constrained RL setting, any two reward functions that belong to the same equivalence class will give rise to the same optimal policy.
\end{lemma}

The proofs are elementary and can be found in Appendix \ref{app:lemma1}. The first lemma highlights the well-established issue of non-uniqueness in the Plackett-Luce model family \cite{plackett1975analysis}. This ambiguity necessitates further identifiability assumptions in order to provide any guarantees for the MLE solutions of Eq. \ref{eq:reward_model} \cite{bong2022generalized}. The second lemma emphasizes that all reward functions within a given equivalence class yield identical optimal policies, thus our primary aim is to recover any representative from the optimal equivalence class. We formally establish the following Theorem in Appendix~\ref{app:thm1}:

\begin{theorem}\label{thm:main}
    Under certain mild conditions, every reward class that is compatible with the Plackett-Luce model (and in particular Bradley-Terry) can be expressed via the following reparameterization: ${r(x, y) = \beta \log \frac{\pi(y\mid x)}{\pi_\text{ref}(y\mid x)}}$ for some model $\pi(y\mid x)$ and a specified reference distribution $\pi_\text{ref}(y \mid x)$.
\end{theorem}

\begin{sproof}
    Take any reward function $r(x, y)$; it defines a corresponding optimal policy $\pi_r(y \mid x)$ given by Eq. \ref{eq:op_policy}. Our goal is to show that a reward function from the same equivalence class as $r$ admits the above reparameterized form. Define the mapping $f$ as
\begin{equation}
    f(r; \pi_\text{ref}, \beta)(x, y) = r(x, y) - \beta\log\sum_{y}\pi_\text{ref}(y\mid x)\exp\left(\frac{1}{\beta}r(x, y)\right)
\end{equation}
The function $f$ adjusts the reward by subtracting the log-partition function evaluated under $\pi_r$, resulting in a normalization dependent only on $x$. Thus, $f(r; \pi_\text{ref}, \beta)(x, y)$ remains in the same equivalence class as $r(x, y)$. Substituting the right-hand side of Eq.~\ref{eq:main_eq} (applicable for any reward function), we observe $f(r; \pi_\text{ref}, \beta)(x, y) = \beta \log \frac{\pi_r(y\mid x)}{\pi_\text{ref}(y\mid x)}$. Consequently, this construction shows that one may represent any member of the equivalence class of $r$ with the desired parametric form, thus incurring no loss of generality when using this reward model reparameterization.
\end{sproof}

Theorem~\ref{thm:main} may alternatively be interpreted as providing a precise characterization of the particular reward function selected by the DPO reparameterization within each equivalence class. Specifically, this is the reward function that fulfills:

\begin{equation}\label{eq:lag_p}
     \sum_{y}\underbrace{\pi_\text{ref}(y\mid x)\exp\left(\frac{1}{\beta}r(x, y)\right)}_{=\pi(y\mid x)\text{, using Thm.~\ref{thm:main} reparam.}} = 1,
\end{equation}

In other words, $\pi(y\mid x)$ constitutes a proper probability distribution, maintaining positivity and normalization. Notably, in reference to Eq.~\ref{eq:op_policy}, Eq.~\ref{eq:lag_p} corresponds to the partition function for the optimal policy derived from the reward function $r(x, y)$. The essential contribution of the DPO algorithm lies in its ability to impose specific conditions on the generally under-constrained Plackett-Luce (as well as Bradley-Terry) family of preference models. These constraints do not alter the set of representable reward models, but rather ensure that the optimal policy in Eq. \ref{eq:op_policy} is rendered analytically solvable for any prompt $x$.

\subsection{Instability of Actor-Critic Algorithms} Our framework also enables the identification and analysis of instability phenomena encountered with standard actor-critic methods in RLHF, such as PPO. In accordance with the RLHF pipeline, we focus on the RL fine-tuning stage as discussed in Section \ref{section:prelims}. A connection can be drawn here to the control as inference perspective \cite{levine2018reinforcement}, as it pertains to the constrained RL formulation in \ref{eq:RL}. Considering a parameterized policy $\pi_{\theta}(y\mid x)$, the goal is to minimize $\mathbb{D}_{\text{KL}}[\pi_{\theta}(y|x) \mid \mid \pi^*(y\mid x)]$, where $\pi^*$ denotes the optimal policy from Eq. \ref{eq:optimum_model} under the reward $r_{\phi}(y, x)$. Algebraic manipulation yields the following optimization target:

\begin{equation}\label{eq:AC}
    \max_{\pi_{\theta}}\mathbb{E}_{\pi_{\theta}(y\mid x)}\bigg[\underbrace{r_{\phi}(x, y) -\beta\log\sum_{y}\pi_\text{ref}(y\mid x)\exp\left(\frac{1}{\beta}r_{\phi}(x, y)\right)}_{f(r_{\phi}, \pi_\text{ref}, \beta)} - \underbrace{\beta\log\frac{\pi_{\theta}(y\mid x)}{\pi_\text{ref}(y\mid x)}}_{\text{KL}}\bigg]
\end{equation}

This objective has also been the focus of optimization in previous studies 
\citep{ziegler2020finetuning, stiennon2022learning, bai2022training, ouyang2022training}, utilizing the DPO-equivalent reward corresponding to the class of $r_{\phi}$. Here, the normalization factor within $f(r_{\phi}, \pi_\text{ref}, \beta)$ can be understood as the soft value function of the reference policy $\pi_\text{ref}$. Although this normalization does not influence the optimal policy itself, its omission can result in policy gradients with substantial variance, potentially destabilizing the learning process. Addressing this normalization can involve introducing a learned value function; however, optimizing such functions is itself challenging. Previous literature has often opted to normalize rewards by employing a baseline derived from human completions, amounting to a single-sample Monte-Carlo approximation for the normalization. By contrast, the DPO reparameterization leads to a reward structure that obviates the need for such baselines.

\section{Experiments}
We empirically investigate how well DPO learns policies from preference data in this section. We start with controlled text-generation experiments to examine how DPO manages the trade-off between reward maximization and minimizing KL-divergence to the reference policy, comparing its efficiency with standard preference-based methods such as PPO. We then proceed to assess DPO’s scalability and efficacy on larger language models and more complex RLHF challenges, such as tasks in summarization and dialogue. Our results indicate that, even with minimal hyperparameter adjustment, DPO typically matches or outperforms strong baselines like PPO-based RLHF and outperforms strategies like selecting the best trajectory from $N$ candidates based on a learned reward. Details on experimental protocols are provided here, with further discussion available in Appendix~\ref{app:exp_details}.

\textbf{Tasks.} In our study, we investigate three open-ended text generation tasks. Across all tasks, the algorithms are tasked with learning a policy using a dataset of preferences $\mathcal{D}=\bigl\{x^{(i)}, y_w^{(i)}, y_l^{(i)}\bigr\}_{i=1}^N$. In the \textbf{controlled sentiment generation} task, $x$ corresponds to a movie review prefix from the IMDb dataset \cite{maas-EtAl:2011:ACL-HLT2011}, and the goal is for the policy to produce $y$ exhibiting positive sentiment. To facilitate a controlled study, we automatically generate preference pairs for each output using a pre-trained sentiment classifier, ensuring that $p(\text{positive}\mid x,y_w)>p(\text{positive}\mid x,y_l)$. For SFT, GPT-2-large is fine-tuned to convergence using the training split of IMDb reviews (see App~\ref{app:sentiment_details} for further specifics). For \textbf{summarization}, $x$ consists of a Reddit forum post, with the policy expected to produce a summary $y$ capturing the main ideas. Consistent with earlier works, we utilize the Reddit TL;DR summarization dataset \citep{volske-etal-2017-tl}, alongside human preference annotations collected by \citeauthor{stiennon2022learning}. The SFT baseline is a model fine-tuned on human-created Reddit post summaries,\footnote{\url{https://huggingface.co/CarperAI/openai_summarize_tldr_sft}} and RLHF training is conducted using the TRLX \citep{leandro_von_werra_2023_7790115} framework. The human preferences were recorded by \citeauthor{stiennon2022learning} based on outputs from a different, yet comparably trained, SFT model. The third task, \textbf{single-turn dialogue}, involves $x$ as a human input, which may range from scientific inquiries to personal advice requests. Here, the policy is expected to generate an informative and relevant reply $y$; data comes from the Anthropic Helpful and Harmless dialogue dataset \citep{bai2022training}, comprising 170,000 conversations between human users and an automated assistant. Each conversation concludes with two candidate responses from a large, unspecified language model, with one labeled as preferred by humans. As there is no available pre-trained SFT model for this dataset, we fine-tune a standard language model solely on the preferred completions to construct the SFT baseline.

\textbf{Evaluation.} We employ two distinct strategies for assessment in our experiments. In the controlled sentiment generation scenario, we focus on how well each algorithm optimizes the constrained reward maximization goal by mapping the trade-off frontier between realized reward and KL-divergence from the reference policy; this is possible because we possess the actual reward function in the form of a sentiment classifier. In contrast, since the true reward function is generally inaccessible in practical settings, we turn to evaluating methods by their \textit{win rate} relative to a baseline policy. Here, we use GPT-4 to approximate human judgment when assessing the quality of summaries or helpfulness of responses in summarization and single-turn dialogue tasks, respectively. For summarization, baseline comparisons are made against the reference summaries from the test split, while for dialogue, the baseline is the test set’s preferred response. Although prior research indicates that large language models (LMs) may serve as superior automated evaluators compared to conventional metrics \citep{Chen2023ExploringTU}, we perform a human study to validate our selection of GPT-4 as the evaluator, as described in Sec.~\ref{sec:human-judgments}. Our analysis reveals a strong alignment between GPT-4 and human judgments, with human-GPT-4 agreement rates matching or exceeding those observed between human annotators.

\textbf{Methods.} Beyond DPO, we benchmark several prevalent methods for training language models according to human preferences. For baseline prompting, we utilize zero-shot prompts with \textbf{GPT-J} \citep{gpt-j} in the summarization experiments and two-shot prompts with \textbf{Pythia-2.8B} \citep{biderman2023pythia} in the dialogue tasks. We additionally consider both the \textbf{SFT} approach and \textbf{Preferred-FT}—the latter being a model subjected to supervised fine-tuning on selected outputs $y_w$, derived from either the SFT model (used in controlled sentiment and summarization settings) or a general-purpose language model (used in single-turn dialogue). Another approach is the \textbf{Unlikelihood} method~\citep{welleck2019neural}, which trains policies to increase the likelihood of $y_w$ while simultaneously \textit{decreasing} the likelihood of $y_l$; here, an optional hyperparameter $\alpha\in[0,1]$ scales the ‘unlikelihood’ loss term. We further assess the \textbf{PPO} algorithm \citep{schulman2017proximal}, leveraging reward models learned from preference data, as well as \textbf{PPO-GT}, which represents an oracle utilizing the ground-truth reward function in the controlled sentiment environment. For sentiment-related experiments, we evaluate both an off-the-shelf PPO-GT implementation \cite{leandro_von_werra_2023_7790115} and an adapted version where rewards are normalized and additional hyperparameter tuning is performed for improved efficacy (these modifications are also applied to PPO with learned rewards). Lastly, we employ the \textbf{Best of $N$} strategy, wherein $N$ candidate completions are sampled from the SFT model (or Preferred-FT for dialogue), and the response yielding the highest score—determined by the reward model from the preference data—is selected. Despite its robust performance, this method is computationally demanding, as it necessitates sampling $N$ outputs per query during evaluation, making it impractical for even moderately large values of $N$.

\subsection{How well can DPO optimize the RLHF objective?}

The KL-constrained reward maximization objective commonly adopted in RLHF algorithms aims to strike a balance between increasing the reward and preventing the policy from straying too far from the reference policy. Consequently, when evaluating different algorithms, it is crucial to consider both the reward obtained and the associated KL divergence; achieving marginally greater reward at the cost of substantially increased KL divergence may not be beneficial. Figure~\ref{fig:frontier-tldr-main} presents the trade-off between reward and KL divergence—the so-called frontier—for several algorithms applied to the sentiment task. For each method, we perform several training runs, varying a hyperparameter that controls policy conservativeness (target KL $\in\{3,6,9,12\}$ for PPO, $\beta \in \{0.05,0.1,1,5\}$, and $\alpha\in\{0.05,0.1,0.5,1\}$ for unlikelihood, as well as different random seeds for preferred-FT), totaling 22 runs. At intervals of 100 training steps, and continuing until convergence, we assess each trained policy on a set of test prompts, measuring both the mean reward (according to the true reward function) and the mean sequence-level KL\footnote{Here, the sequence-level KL is calculated as the aggregate of per-timestep KL-divergences.} with respect to the reference policy, $\text{KL}\left(\pi\mid \mid \pi_\text{ref}\right)$. Our experiments reveal that DPO produces the most favorable reward-KL frontier: it attains the highest reward for any given KL value, and does so while keeping KL low. This finding is especially noteworthy for several reasons. Firstly, although both DPO and PPO seek to optimize the same objective, DPO does so with substantially greater efficiency; DPO's reward/KL curve consistently outperforms that of PPO. Secondly, DPO surpasses PPO in this trade-off even when PPO is supplied with exact reward information (PPO-GT).

\subsection{Can DPO scale to real preference datasets?}
\label{sec:dpo-real-datasets}
We proceed by assessing how well DPO performs when fine-tuned for tasks involving summarization and single-turn dialogue. In the case of summarization, automatic metrics like ROUGE have been shown to align poorly with human judgments~\citep{stiennon2022learning}, motivating previous approaches that leverage PPO-based fine-tuning with human preferences to generate more useful summaries. Our evaluation compares several approaches by generating summaries for the test partition of the TL;DR dataset, then calculating the mean win rate over the reference summaries in the test set. For all evaluated methods, completions are generated across a range of temperatures from 0.0 to 1.0, with corresponding win rates presented in Figure~\ref{fig:frontier-tldr-main} (right). All methods—DPO, PPO, and Preferred-FT—are initialized from the same GPT-J SFT model\footnote{\url{https://huggingface.co/CarperAI/openai_summarize_tldr_sft}}. In our experiments, DPO achieves a win rate close to 61\% at temperature 0.0, outperforming PPO, which attains about 57\% at its optimal setting. Additionally, DPO reaches a greater peak win rate than the best among $N$ baseline approaches. It is worth noting that we did not perform extensive optimization of DPO's $\beta$ parameter, suggesting the reported results might be a conservative estimate of its capabilities. Notably, DPO demonstrates greater stability with respect to temperature variation than PPO, whose win rates can fall to match those of the original GPT-J model at higher temperatures. Preferred-FT, in contrast, shows minimal improvement relative to the SFT baseline. We further conduct a direct comparison between DPO and PPO via human evaluation in Section~\ref{sec:human-judgments}; there, samples from DPO at temperature 0.25 are chosen over PPO samples at the same temperature in 58\% of cases.

For single-turn dialogue tasks, we assess the various approaches on a portion of the test set from the Anthropic HH dataset \citep{bai2022training} that consists of interactions involving just one human-assistant exchange. To measure win rates across methods, GPT-4 is employed to compare candidate completions with the preferred responses in the test set, treating the latter as the ground truth. Since no widely-adopted SFT model is established for this benchmark, our procedure begins with a pre-trained Pythia-2.8B. We apply Preferred-FT, fine-tuning a reference model on selected completions to align its output distribution accordingly, after which we proceed to train with DPO. For further comparison, we include both the strongest completion out of 128 Preferred-FT samples—an empirically sufficient number for this benchmark, as indicated by the plateau shown in Appendix Figure~\ref{fig:best-of-n}—as well as a 2-shot prompt strategy using the base Pythia-2.8B model. Our findings indicate that DPO performs as well as, or outperforms, competing approaches across optimal temperature settings. We also evaluate an RLHF system trained via PPO on the Anthropic HH data \footnote{\url{https://huggingface.co/reciprocate/ppo_hh_pythia-6B}}, implemented using a reputable framework \footnote{\url{https://github.com/CarperAI/trlx/tree/main/examples/hh}}, but are unable to identify a prompt or temperature configuration that surpasses the performance of the unaltered Pythia-2.8B. Drawing on results from TL;DR and recognizing that both methods maximize the same reward signal, we use Best of 128 as an approximate surrogate for PPO-based performance. In summary, DPO stands out as the only approach that both improves over the original preferred completions from the Anthropic HH dataset and does so with far less computational cost compared to the Best of 128 reference. Figure~\ref{fig:dialogue-main} additionally demonstrates that DPO attains peak performance within a relatively small number of steps.

\subsection{Generalization to a new input distribution}

To more thoroughly assess the robustness of PPO and DPO to distributional changes, we re-examine the policies learned from our Reddit TL;DR summarization task on an alternative dataset: the test set of news articles from CNN/DailyMail \citep{nallapati-etal-2016-abstractive}. For this cross-domain evaluation, we apply the best-performing sampling temperatures from the TL;DR task (0 and 0.25). The findings are shown in Table~\ref{tab:ood}. We calculate the win rate of each method compared to ground-truth summaries, employing the same GPT-4 (C) prompt used previously on Reddit TL;DR data, with the modification that “forum post” is replaced by “news article.” Within this new distribution, DPO maintains a strong advantage over PPO. These results serve as preliminary evidence that DPO can match or exceed the generalization abilities of PPO, even though DPO is trained without access to the additional unlabeled Reddit TL;DR prompts incorporated by PPO.

\subsection{Validating GPT-4 judgments with human judgments}
\label{sec:human-judgments}
We implement a human evaluation to assess the trustworthiness of GPT-4’s assessments, using the outputs from our TL;DR summarization experiment and two distinct GPT-4 prompt formulations. The \textbf{GPT-4 (S)} (simple) prompt requests a comparison based solely on which summary better captures the key information. In contrast, the \textbf{GPT-4 (C)} (concise) prompt directs attention not only to informativeness but also to conciseness; we use this prompt as initial observations indicate GPT-4 (under the simple prompt) tends to prefer lengthier, more repetitive responses than human evaluators. The complete wording of each prompt can be found in Appendix~\ref{app:prompts}. To achieve a representative comparison of model outputs, we select three sampling strategies: the highest performing (DPO, temp. 0.25), lowest performing (PPO, temp. 1.0), and an intermediate case (SFT, temp. 0.25), each contrasted against PPO’s best-performing greedy decoding output. Across both prompts, agreement rates between GPT-4 and human raters closely track human inter-annotator agreement, suggesting that GPT-4 may function as an adequate substitute for human evaluation (note that, due to a limited number of raters, we gather multiple human judgments only for the DPO and PPO-1 comparisons). Notably, the \textbf{GPT-4 (C)} prompt yields win rates that are generally more aligned with human judgments; accordingly, this prompt is adopted for primary results in Section~\ref{sec:dpo-real-datasets}. For further methodological specifics, including a description of the user interface and the roster of participating volunteers, consult Appendix~\ref{app:human-study}.

\section{Discussion}
Preference-based learning offers a robust and flexible approach for developing proficient and aligned language models. In this work, we have presented DPO, a straightforward method for leveraging preferences in language model training that does not rely on reinforcement learning. Instead of adapting the preference learning process to fit conventional RL frameworks solely to utilize pre-existing RL methodologies, DPO constructs an explicit connection between language model policies and corresponding reward functions. This approach allows direct optimization of language models to adhere to human preferences via a basic cross-entropy objective, circumventing the need for reinforcement learning while preserving general applicability. Remarkably, DPO attains performance comparable to or better than established RLHF techniques, including those utilizing PPO, with minimal hyperparameter adjustment; thus, DPO significantly lowers the practical barriers for training language models grounded in human preference data.

\textbf{Limitations \& Future Work.} Our findings prompt several avenues for continued research. It is important to understand the extent to which policies trained via DPO exhibit out-of-distribution generalization compared to models trained on explicit reward signals. Preliminary evidence suggests DPO yields generalization performance on par with PPO-driven alternatives, though a broader evaluation is warranted. Another open question concerns the potential of utilizing self-labeled data through the DPO policy for efficiently leveraging unlabeled prompts. Additionally, the dynamics of reward over-optimization within direct preference optimization frameworks, and whether the slight dip in outcomes depicted in Figure~\ref{fig:dialogue-main}-right illustrates this issue, merit further investigation. While our current study encompasses models of up to 6B parameters, examining the scalability of DPO for much larger, state-of-the-art architectures represents a compelling research trajectory. Concerning evaluation methodology, we observe that GPT-4 derived win rates are sensitive to prompt formulation; future research could focus on identifying optimal practices for extracting reliable automated evaluations. Beyond aligning language models, the utility of DPO potentially extends to training generative models in domains outside of natural language, signaling diverse possibilities for future exploration.

\end{document}